% Section 9.1, from lectures/L09/appendix/a_the_pair.md.

\section{The pair, the tail, and the two factors of two}
\label{c9:app:a}

The first stage in this book whose purpose is to ignore something.

\subsection{Two halves and one current}
\label{c9:sec:a1}
\bookfigure{lectures/L09/appendix/images/differential_pair.png}{A differential pair: two matched NPN transistors with their emitters joined to a tail current source, collector resistors to the positive rail, both bases driven, and the output taken at one collector.}{c9:fig:differentialpair}

Two matched transistors with their emitters tied together and a \textbf{tail} current source below. The tail fixes the total, $I_{C1} + I_{C2} = I_{tail}$; the inputs decide only how it is \textbf{divided}. Both properties of the circuit follow from that:

\begin{itemize}
  \item \textbf{A differential input redistributes the current} between the halves. The tail node barely moves, because the sum is unchanged.
  \item \textbf{A common-mode input moves both halves together,} which the tail cannot supply, so the tail node has to move instead.
\end{itemize}
\textbf{Each side runs at half the tail current.} For a 2 mA tail, 1 mA per side:

\[
  r_e = \frac{V_T}{I_{tail}/2} = 26\ \Omega
\]

\textbf{26 ohm, not 13.} Reading $r_e$ from the tail current rather than the side current is the commonest slip in this chapter, and it is a factor of two on top of the two below.

\subsection{Differential gain, and both factors of two}
\label{c9:sec:a2}
Apply $+v_d/2$ to one base and $-v_d/2$ to the other. The tail node does not move, so each half is a common-emitter stage with its emitter at signal ground: \textbf{exactly Chapter~\ref{c7:ch:smallsignal}'s circuit}, no new derivation.

\[
  A_1 = -\frac{R_C}{r_e} \cdot \frac{v_d/2}{v_d} = -\frac{R_C}{2 r_e}
\]

For 10 kilohm and a 2 mA tail, $-192$.

\begin{itemize}
  \item \textbf{The first two is the input split.} A differential input of $v_d$ puts only $v_d/2$ on each base. Unavoidable: it is what differential means.
  \item \textbf{The second two is the output.} Take the difference between the two collectors and the gain is $-R_C/r_e = -385$, the collectors moving oppositely. Take \textbf{one} collector and the other half is thrown away.
\end{itemize}

\textbf{The second two is recoverable and the first is not.} Recovering it is what the current mirror of \secref{c9:sec:b4} does, and one of the two reasons that mirror exists. This book takes the \textbf{single-ended} output throughout, because that is what feeds the next stage in an operational amplifier.

\subsection{Input and output resistance}
\label{c9:sec:a3}
\textbf{Looking into either base,} Chapter~\ref{c7:ch:smallsignal}'s result with the same factor of two:

\[
  Z_{in(diff)} = 2 h_{FE} r_e
\]

2.6 kilohm at a 2 mA tail with $\beta = 50$: the two bases in series as far as a differential source is concerned. Low, and the reason \secref{c9:sec:b6} prefers MOSFETs here.

\textbf{Looking into a collector,} Chapter~\ref{c7:ch:smallsignal}'s \code{resistanceIntoCollector} with the tail as the degeneration resistor. With a current-source tail that is enormous, so the output resistance is $R_C$, or the mirror's $r_o$ when there is a mirror.

\textbf{And the base current is not zero,} which matters more here than anywhere else. Each base draws $I_{tail}/2\beta$, 20 microamps at a 2 mA tail. Those currents flow through whatever drives the inputs, and \textbf{if the two source resistances differ, the difference becomes an input voltage the amplifier cannot distinguish from signal.}

\subsection{What the small-signal model cannot see}
\label{c9:sec:a4}
Two exponentials dividing one current is a hyperbolic tangent:

\[
  I_{C1} - I_{C2} = I_{tail}\tanh\!\left(\frac{v_d}{2 V_T}\right)
\]

\bookfigure{lectures/L09/appendix/images/diffpair_transfer.png}{The difference between the two collector currents against differential input, an S-shaped curve saturating at plus and minus the tail current, with the small-signal tangent drawn through the origin and the region where the two agree to one per cent shaded.}{c9:fig:diffpairtransfer}

\textbf{The pair is linear over nine millivolts.} Beyond $\pm 9.1$ mV the tanh has fallen 1 per cent below its tangent, and that is the pair's honest input range.

\textbf{And that figure does not depend on the tail current at all.} The tail scales the whole curve and cancels out of the ratio. \textbf{Biasing the pair harder buys gain and no linearity whatever,} which is not what anyone expects and is worth checking in the code.

\textbf{The pair hard-limits.} At $\pm 100$ mV, 96 per cent of the tail has moved to one side and the other is off. Past that the output stops responding: a ceiling, not soft compression.

\textbf{That ceiling is where slew rate comes from.} The pair drives a compensation capacitor, and once fully switched the largest current it can deliver is $I_{tail}$, so the output ramps at $I_{tail}/C$ and no feedback loop can make it faster. Slew rate is not a small-signal parameter, and this is why.

\subsection{What this section is blind to}
\label{c9:sec:a5}
\begin{itemize}
  \item \textbf{Mismatch.} Everything here assumes the two halves are identical. They are not, and \secref{c9:sec:b6} is about what that costs.
  \item \textbf{The tail's own behaviour.} Ideal here; the whole of \secref{c9:app:b} is about it not being.
  \item \textbf{Frequency.} The pair has any stage's Miller problem, mitigated by each half seeing only half the signal. Not treated.
  \item \textbf{Noise}, where a differential input stage earns much of its reputation, and which this book does not cover anywhere.
\end{itemize}
